% Generic fermionic Gaussian pairing combinatorics before QED specialization.
自由高斯理论的高阶关联函数由二点函数决定。玻色场交换两个算符不会改变符号；费米场每交换两个 Grassmann 奇算符都会产生一个负号。因此一般高阶费米关联函数的每一种完全配对都必须带上把原有算符顺序重排到该配对顺序所需的置换宇称。后面的全同费米子交换号与闭合费米子圈负号都来自这一同一代数结构。


设按原始时间序列排列的 $2n$ 个自由费米型场算符统一记为
\[
\Psi_1,\Psi_2,\ldots,\Psi_{2n},
\]
其中每个 $\Psi_r$ 可以是 $\psi$ 或 $\bar\psi$，并把二点缩并记为
\begin{equation}
 C_{rs}\equiv
 \langle0|\Tord\{\Psi_r\Psi_s\}|0\rangle_{\rm contraction}.
 \label{eq:generic-fermion-contraction-dp8}
\end{equation}
只有场类型与量子数允许的二点缩并才非零。令 $\mathfrak P_{2n}$ 表示集合 $\{1,\ldots,2n\}$ 的全部完全配对；对每个配对 $P$，$(-1)^{\pi(P)}$ 定义为把原始算符序列重排成各配对相邻时所产生的费米置换符号。自由高斯费米场的一般 $2n$ 点函数于是为
\begin{equation}
 \boxed{
 \langle0|\Tord\{\Psi_1\Psi_2\cdots\Psi_{2n}\}|0\rangle
 =\sum_{P\in\mathfrak P_{2n}}
 (-1)^{\pi(P)}
 \prod_{(r,s)\in P} C_{rs}.}
 \label{eq:fermion-wick-general-dp73}
\end{equation}
这一“任意高阶高斯关联函数等于所有二点配对之和，并由费米置换宇称定号”的结果通常称为费米 Wick 定理。其母结构是\eqnrefs{eq:fermion-wick-general-dp73}，低阶公式只是代入后的检验。

例如取 $n=2$ 且算符次序为 $\psi_1\bar\psi_2\psi_3\bar\psi_4$，两个允许的完全收缩具有相反宇称，因此
\begin{equation}
 \boxed{
 \langle0|\Tord\{\psi_1\bar\psi_2\psi_3\bar\psi_4\}|0\rangle
 =S_{12}S_{34}-S_{14}S_{32},}
 \label{eq:fermion-wick-four-r68}
\end{equation}
其中 $S_{ij}\equiv\langle0|\Tord\{\psi_i\bar\psi_j\}|0\rangle$。

在微扰展开中，若一组费米缩并首尾相接形成内部闭环，把该闭环上的 Grassmann 场重排为一个有序矩阵迹时会比对应的玻色配对多一次奇置换，因此每个独立闭合费米子圈带有
\begin{equation}
 \boxed{\text{each closed fermion loop contributes an overall factor }-1.}
 \label{eq:generic-closed-fermion-loop-r68}
\end{equation}
具体相互作用顶角只决定环中的矩阵、内部量子数与传播子；这个整体负号在进入 QED、QCD 等具体理论之前已经由费米反对易代数固定。
